\newcommand{\Int}{\mathbb{Z}}
\subsection{Критерий достижимости точки через $D_{n,k}$}
\temporarytext{Перейти к источнику \cite{Leaper_Graphs_Knuth} и ссылаться на приведённый там
результат вместо его повторного доказательства.}

$D_{n,k}(a, b)$ верно тогда и только тогда, когда $a = b + (\pm n,\pm k)$ или $a = b + (\pm k,\pm n)$, где знаки выбираются независимо.

Будем называть точку $b$ \textit{достижимой} из $a$ за $N$ шагов $D_{n,k}$, если существуют такие $b_1, b_2, \dots, b_N$, что:

$$D_{n,k}(a, b_1) \wedge (\bigwedge \limits_{i = 1}^{N - 1} D_{n,k}(b_i, b_{i + 1})) \wedge D_{n,k}(b_N, b)$$

\begin{theorem}
    Пусть $n$ и $k$ --- положительные целые числа и $d = \gcd(n,k),\  p = \frac{n}{d},\ q = \frac{k}{d}$.

    Точка $(x,y)$ достижима из $(0,0)$ через конечное число шагов $D_n,k$ тогда и только тогда, когда выполняются следующие условия:

    \begin{enumerate}
        \item $d \mid x$;
        \item $d \mid y$;
        \item $(2 \nmid p \wedge 2 \nmid q) \Longrightarrow \frac{x}{d} \equiv \frac{y}{d}\ \pmod{2}$.
    \end{enumerate}
\end{theorem}

Критерий связности графа обобщённого шахматного коня является известным
результатом \cite{Connectivity_Knight_Graphs_Rhodes_Wilson}; ниже приведено его доказательство в используемых
обозначениях вместе с описанием отдельных классов достижимости.

\textbf{Доказательство.}

Для доказательства утверждения воспользуемся леммой Безу:

\begin{namedLemma}[лемма Безу]
    $\forall n,k \in \mathbb{Z}\ \exists a, b \in \mathbb{Z}\ |\ \gcd(n, k) = an + bk$.
\end{namedLemma}

\textbf{Доказательство необходимости.}

По определению достижимости существует цепочка $(x_1, y_1), (x_2, y_2), \dots, (x_N, y_N)$ такая, что:

$$D_{n,k}((0,0), (x_1, y_1)) \wedge (\bigwedge \limits_{i = 1}^{N - 1} D_{n,k}((x_i, y_i), (x_{i + 1}, y_{i + 1}))) \wedge D_{n,k}({x_N, y_N}, (x,y))$$.

Заметим, что

$$D_n,k((x_1, y_1), (x_2, y_2)) \Leftrightarrow x_1 = x_2 \pm n\ \wedge\ y_1 = y_2 \pm k\ \vee\ x_1 = x_2 \pm k\ \wedge\ y_1 = y_2 \pm n \Rightarrow$$ 

$$\Rightarrow x_1 = \alpha \cdot d + x_2\ \wedge\ y_1 = \beta \cdot d + y_2$$

откуда $d \mid x_1 - x_2,\ d \mid y_1 - y_2$. Следовательно, $d \mid x\ \wedge\ d \mid y$.
\ms

Теперь предположим, что $p$ и $q$ оба нечётные. После деления координат
каждого шага на $d$ получается один из векторов

$$(\pm p,\pm q)\quad\text{или}\quad(\pm q,\pm p).$$

Обе координаты такого вектора нечётные, следовательно каждый шаг одновременно
меняет чётность обеих координат. В точке $(0,0)$ координаты имеют
одинаковую чётность, следовательно, это остаётся верным во всех
достижимых точках:

$$\frac{x}{d}\equiv\frac{y}{d}\pmod 2.$$

Необходимость доказана. \hfill $\square$
\bs

\textbf{Доказательство достаточности.}

Назовем \textit{разрешенными} сдвиги $(a, b)$, для которых из достижимости $(x,y)$ следует достижимость $(x + a, y + b)$.

Пусть сдвиг $(a, b)$ разрешен. Тогда выполнены следующие свойства:
\begin{enumerate}
    \item $(-a, b)$ разрешен.
    \item $(b, a)$ разрешен.
    \item Если сдвиг $(a_1, b_1)$ разрешен, то $(a \pm a_1, b \pm b_1)$ также разрешен.
\end{enumerate}

По свойству 1 сдвиги $(2n, 0)$ и $(2k, 0)$ разрешены:
$$(n,k) + (n,-k) = (2n,0)$$
$$(k,n) + (k,-n) = (2k,0)$$

Покажем, что тогда $(2d,0)$ разрешен.

По лемме Безу существуют целые $a,b$, для которых $d = an + bk$. Тогда: 

$$|a|(\text{sign}(a) \cdot 2n,0) + |b|(\text{sign}(b) \cdot 2k,0) = (2an+2bk,0) = (2d,0)$$

Значит, $(2d,0)$ разрешен. Аналогично разрешен $(0,2d)$.
\ms

Таким образом, разрешены все сдвиги вида $(2d \cdot i,2d\cdot j)$. Рассмотрим два случая: 
\ms

\textbf{Случай 1: одно из чисел $p$, $q$ чётное.}
\ms

Так как $p$ и $q$ взаимно просты, второе число в этом случае нечётно. Не умаляя общности $p$ нечётно, а $q$ чётно. 

Так как $p$ нечётно, а $q$ чётно, числа $(p-1)/2$ и $q/2$ целые. 

Тогда разрешен сдвиг $$(n,k) - \frac{p-1}{2}(2d,0) - \frac{q}{2}(0,2d) = (dp,dq) - (d(p-1),0) - (0,dq) = (d,0)$$

Следовательно, $(d,0)$ разрешен. Аналогично, $(0,d)$ разрешен. Тогда разрешены все сдвиги вида $(d \cdot i,d\cdot j)$. 
\ms

Отсюда $(x, y) = (d \cdot i, d\cdot j)$ достижима из (0,0).
\bs

\textbf{Случай 2: $p$ и $q$ нечётны.}

Пусть

$$x = du,\ y = dv$$

где $u$ и $v$ имеют одинаковую чётность (по условию 3 теоремы).

Если $u$ и $v$ чётны, числа $u/2$ и $v/2$ целые. Поэтому

$$(x,y) = \frac{u}{2}(2d,0) + \frac{v}{2}(0,2d)$$

Если $u$ и $v$ нечётны, числа $u-p$ и $v-q$ чётны. Поэтому числа $(u-p)/2$ и $(v-q)/2$ целые. Тогда:

$$(n,k) + \frac{u-p}{2}(2d,0) + \frac{v-q}{2}(0,2d) = (dp,dq) + (d(u-p),0) + (0,d(v-q)) = (du,dv) = (x,y),$$

то есть $(x,y)$ достижима из $(0,0)$.
\bs

Таким образом, указанные условия всегда достаточны. \hfill $\square$
\bigskip

\begin{corollary}[\textbf{Классы достижимости}]
Точки $(x_1,y_1)$ и $(x_2,y_2)$ находятся в одном классе достижимости тогда и только тогда, когда выполнено одно из следующих условий:
\begin{enumerate}
    \item Если одно из чисел $p$, $q$ чётно, то
    $$x_1\equiv x_2\pmod d \ \land\ y_1\equiv y_2\pmod d.$$
    \item Если $p$ и $q$ нечётны, то
    $$\exists i\in\mathbb Z:\quad x_1\equiv x_2+i\cdot d\pmod{2d}
    \ \land\ y_1\equiv y_2+i\cdot d\pmod{2d}.$$
\end{enumerate}
\end{corollary}
\bs

Обозначим класс достижимости относительно $D_{n,k}$, содержащий точку $(x,y)$, как $RC_{n,k}(x,y)$.

\subsubsection{Эквивалентность бинарных отношений $D_{n,k}$}

Определим два подсемейства $D_{n,k}$:
\begin{itemize}
    \item $S_d \leftrightharpoons D_{d,0}$;
    \item $D_d \leftrightharpoons D_{d,d}$.
\end{itemize}


\begin{lemma}
    Пусть $p,q$ --- положительные целые числа и $\gcd(n,k) = 1$, $n \equiv k \pmod 2$. Тогда $D_{n,k} \approx S_1$.
\end{lemma}
\textbf{Доказательство.} 

Докажем, что множество автоморфизмов $\langle\Int^2, D_{n,k}\rangle$ сохраняет $\langle\Int^2, S_1\rangle$.

Определим $S_{n,k} = \{(\pm n,\pm k),(\pm k, \pm n)\}$. Тогда выполнено $D_{n,k}(x,y) \Iff x - y \in S_{n,k}$.

Тогда структуру $\langle\Int^2, D_{n,k}\rangle$ можно рассматривать как граф Кэли $\Cay(\Int^2, S{n,k})$. Воспользуемся теоремой~1.4 из источника~\cite{Cayley_Graphs_Nilpotent_Groups_Morris_Morris_Verret}, приведённой ниже на языке 
оригинала.\footnote{Доказательство используемого нами случая для $\Int^2$ опирается на работу A.A. Рябченко~\cite{Cayley_Graphs_Free_Abelian_Group_Ryabchenko}.}

\begin{theorem}
    \label{theorem:Cayley_Graphs_Isomorphisms}
    Assume:
    \begin{itemize}
        \item $G_1$ and $G_2$ are torsion-free, nilpotent groups.
        \item $S_i$ is a finite, symmetric generating set of $G_i$, for $i=1,2$.
    \end{itemize}
    Then every isomorphism from $\operatorname{Cay}(G_1;S_1)$ to 
    $\operatorname{Cay}(G_2;S_2)$ is an affine bijection.
\end{theorem}

Пусть $G_1 = G_2 = \Int^2$, $S_1 = S_2 = S_{n,k}$. Тогда по теореме \ref{theorem:Cayley_Graphs_Isomorphisms} выполнено\footnote{Строго говоря, из теоремы следует утверждение только в одну сторону, в обратную сторону утверждение следует из определения автоморфизма.}:
$$\varphi\in\operatorname{Aut}\bigl(\operatorname{Cay}(\Int^2;S_{n,k})\bigr)
  \Iff
  \exists t\in\mathbb Z^2\
  \exists A\in\operatorname{GL}(2,\mathbb Z):
  \left(A(S_{n,k})=S_{n,k} \land \varphi(x)=Ax+t\right).$$

Из условия $A(S_{n,k})=S_{n,k}$ следует, что $A = 
\begin{pmatrix}
  \pm 1 & 0\\
  0 & \pm 1
\end{pmatrix} 
\vee 
A = \begin{pmatrix}
  0 & \pm 1\\
  \pm 1 & 0
\end{pmatrix}.$

Откуда любое преобразование вида $Ax + t$ сохраняет отношение $D_{1, 0}$ на $\Int^2$.

Рассмотрим полное вверх расширение $\Int^2$. Согласно \cite{Upward_Complete_Structures_Semenov_Soprunov}, $(\Int^2)^*$ представляет собой счетное количество решеток $\Int^2$. 
Автоморфизмы $(\Int^2)^*$ в таком случае будут иметь вид перестановок галактик между собой и применения к каждой из них автоморфизма, описанного выше. Такие автоморфизмы также сохраняют $D_{1, 0}$.

Воспользуемся теоремой Свенониуса\footnote{Текст теоремы взят из источника \cite{Svenonius_Theorem_Model_Theory_Hodges} и приведен в оригинале.}:

\begin{namedTheorem}[Теорема Свенониуса]
    Let $M=\langle A,\Sigma\rangle$ be a countable structure, and $\Sigma \subset \Sigma'$. Then, for any relation $P \in \Sigma$, the following two conditions are equivalent:
    \begin{enumerate}
        \item $P$ is not definable in $\Sigma'$.
        \item There exists an elementary extension $M' = \langle A', \Sigma' \rangle$ of the structure $M'$ and a permutation $\varphi$ of the domain of $A'$, preserving all relations in $\Sigma'$ and not preserving $P$.
    \end{enumerate}
\end{namedTheorem}

Поскольку условие 2 неверно для $M = \langle \Int^2, D_{n,k} \rangle$ и $P = D_{1, 0}$:
$$D_{n,k} \succ D_{1, 0}$$

По определению:

$$E = D_{1, 0} \succ D_{n,k}$$

Тогда $D_{1, 0} \approx D_{n,k} \Then S_1 \approx D_{n,k}$. \hfill $\square$


\bs
\begin{theorem}
    Пусть $n,k$ --- положительные целые числа, $d=\gcd(n,k)$, $p=\frac{n}{d}$ и $q=\frac{k}{d}$. Тогда:
    
    $$
        D_{n,k} \approx
        \begin{cases}
            DS_d, & p\not\equiv q\pmod 2,\\
            DD_d, & p\equiv q\pmod 2.
        \end{cases}
    $$
\end{theorem}
\textbf{Доказательство.} 
Рассмотрим два случая:

\textbf{Случай 1: $p\not\equiv q\pmod 2$.} 

\bs
\textbf{Случай 2: $p\equiv q\pmod 2$.}
